\section{Selection of paths at several exponential ends}\label{sec:path}

The reduction used to prove Theorem~\ref{thm:tree-inequality} needs a path on
which every original parameter has finite weighted variation.  The parameters
include several independently diverging neck lengths.  Ordinary analytic
curve selection does not directly apply at those ends.  We give the precise
asymptotic class and the selection argument used below.  In particular, the
argument includes the specialization of several lengths; one-variable
quasianalyticity alone would not supply that step.

\subsection{The class of functions and the selection statement}

Put
\[
 Q(A,c)=\{L\in\C:\Re L>A,\ |\Im L|<c(\Re L)^2\}.
\]
An ordinary parameter is a variable in a fixed complex neighborhood of a
finite real point.  A real holomorphic function is one satisfying
\(P(\overline L,\overline z)=\overline{P(L,z)}\).

\begin{definition}\label{path:def-admissible}
An admissible family in \(m\) lengths consists of real holomorphic functions
\(P(L_1,\ldots,L_m,z)\), with the following properties.  These properties
are required for every recursively extracted coefficient as well as for the
initial functions.
\begin{enumerate}[label=\textup{(\roman*)}]
\item The functions are bounded on a product of regions \(Q(A,c)\), locally
uniformly in the ordinary parameters.
\item For each length there is an expansion
\begin{equation}\label{path:eq-expansion}
 P(L,z)\sim\sum_{d\in D_i}e^{-dL_i}
             \sum_{k=0}^{k_d}L_i^k P_{dk}(L_{\widehat i},z).
\end{equation}
Here \(D_i\subset[0,\infty)\) is locally finite, each \(k_d\) is finite,
and \(k_0=0\) if \(0\in D_i\).  For every cutoff \(B\), subtracting the
terms with \(d\le B\) leaves a remainder bounded by
\(C_Be^{-(B+\epsilon_B)\Re L_i}\), for some \(\epsilon_B>0\), on
possibly smaller quadratic regions.  Bounds are uniform on each fixed
smaller ordinary-parameter polydisc.
\item For every subset of expanded lengths, all its coefficient formulas,
at all orders, are holomorphic on one common tail band
\[
 \prod_{j\text{ unexpanded}}
 \{\Re L_j>A',\ |\Im L_j|<c'\},\qquad c'>0,
\]
times one common open ordinary-parameter neighborhood.  These domains may
be shrunk once when a finite number of families are used together.  The
constants and thresholds in asymptotic \emph{estimates} may depend on the
order; the common domains of \emph{analyticity} do not.
\item If a subset of real parts has the form
\(\Re L_i=c_iU+o(U)\), where \(c_i>0\), the corresponding iterated
expansion can be grouped by weights \(\sum_i c_id_i\).  The grouped
weights are locally finite, and each finite cutoff has a remainder of
strictly better exponential order in \(U\).  The assertion holds on the
complex neighborhoods in \textup{(i)} and locally uniformly on the common
ordinary-parameter neighborhood.  It is enough to establish it by
successive single-length expansions.
\end{enumerate}
When \(m=0\), admissible families are ordinary holomorphic germs.
\end{definition}

The distinction in \textup{(iii)} will be used when an expansion has
infinitely many zero coefficients.  It is not an assertion that all
remainder estimates hold beyond a single order-independent threshold.

\begin{theorem}[Path selection]\label{thm:path-selection}
Let \(P_1,\ldots,P_s\) be an admissible finite family, and let a finite
Boolean combination of conditions \(P_j>0\), \(P_j=0\), and \(P_j<0\)
hold along a sequence
\[
 L_i^{(n)}\longrightarrow+\infty,\qquad z^{(n)}\longrightarrow z_0.
\]
There is a real analytic map \(u\mapsto(L(u),z(u))\), defined for all
sufficiently large real \(u\), with the same limits and satisfying those
conditions on its tail.  Each actual coordinate \(L_i(u)\) and
\(z_a(u)\) is eventually monotone or constant.  Given any additional
finite admissible list, its values on the path can also be required to be
eventually monotone or constant.  Finite prescribed limits of functions
can be imposed at the same time.
\end{theorem}

A path germ here means analyticity at every sufficiently large finite
\(u\); analyticity at the limiting endpoint is not asserted.  To impose
\(P\to\ell\), introduce an ordinary variable \(v\to\ell\) and impose
\(P-v=0\).  If there are no length variables, add an unused length.  This
also provides positive infinitesimals for strict sign conditions at an
ordinary limiting point.

\begin{corollary}[The variation consequence]\label{path:variation}
Suppose admissible functions \(F,J\), with \(J\ge0\), satisfy
\[
 |dF|\le C\sqrt J\left(\sum_a|dz_a|
                  +\sum_i e^{-cL_i}|dL_i|\right),\qquad c>0,
\]
on real parameter domains.  There cannot be a sequence with
\(L_i\to\infty\), an interior ordinary-parameter limit, and
\(F\ne0\), \(F^2\ge c_0J\), \(F\to0\), where \(c_0>0\).
\end{corollary}
\begin{proof}
Theorem~\ref{thm:path-selection} realizes all these conditions on one
path with monotone actual coordinates.  Along it,
\[
 |d\log|F||\le\frac C{\sqrt{c_0}}
       \left(\sum_a|dz_a|+\sum_i e^{-cL_i}|dL_i|\right).
\]
Each \(z_a\) has finite total variation because it is monotone with
a finite limit.  Each length tends monotonically to infinity on its
tail, and
\[
 \int_{u_0}^\infty e^{-cL_i(u)}|L_i'(u)|\,du
       =c^{-1}e^{-cL_i(u_0)}<\infty.
\]
Thus \(\log|F|\) has finite total variation and a finite limit, whereas
\(F\to0\) forces \(\log|F|\to-\infty\), a contradiction.
\end{proof}

The remaining subsections prove Theorem~\ref{thm:path-selection}.
The geometric construction resumes in Section~\ref{sec:family} and
uses Definition~\ref{path:def-admissible} and
Corollary~\ref{path:variation} as its analytic interface.

\subsection{Sectors, uniqueness, and clocks}

\begin{definition}\label{path:def-sector}
An admissible angular loss is a positive continuous function \(\omega(r)\), for large
\(r\), such that
\[
 \omega(r)\ge(\log r)^{-2},\qquad
 \sum_{n\ge n_0}\sup_{2^n<r<2^{n+2}}\omega(r)<\infty.
\]
Its sector is the open set
\(S_\omega=\{z:|z|>R,\ |\arg z|<\pi/2-\omega(|z|)\}\).
Every positive loss satisfying the displayed bounds has a continuous
admissible enlargement: at dyadic radii choose a summable majorant of
the neighboring shell suprema and interpolate linearly in \(\log r\).
Thus enlarging a loss to make its sector open only shrinks that sector.
Increasing a loss by a fixed factor or taking a continuous majorant of
the suprema over finitely many neighboring dyadic shells is permitted.  Given finitely many losses we
can dominate all of them by another admissible loss.  We can also arrange
that the domination ratio tends to infinity: if \(a_n\ge0\) is summable,
choose increasing constants \(c_n\to\infty\), sufficiently slowly that
\(\sum c_na_n<\infty\).
\end{definition}

The principle that a bounded holomorphic function is determined by its
exponentially accurate expansion is central to Ilyashenko
quasianalyticity; see~\cite[Fact~1.5]{SpeisseggerQA} in the published
version. We prove the following variant with a summable angular loss,
which is the form preserved by the clock changes below.

\begin{lemma}[Sector uniqueness]\label{path:sector-uniqueness}
Suppose \(f\) is bounded and holomorphic on an admissible sector and, on
the central subsector \(|\arg z|\le\pi/4\), satisfies
\(f(z)=O_B(e^{-B|z|})\) for every \(B>0\).  Then \(f=0\).
\end{lemma}
\begin{proof}
Let \(w_n\) be the supremum of the prescribed loss on a fixed enlargement
of the \(n\)-th dyadic shell.  Choose summable \(a_n\ge0\) with
\(a_n/w_n\to\infty\); enlarge them by a fixed factor if necessary.
On the right half-plane define
\[
 A(z)=\sum_n\frac{a_n}{1+z/2^n},\qquad \Phi(z)=z\exp A(z).
\]
The series is normally convergent there, \(A(z)\to0\) uniformly as
\(|z|\to\infty\), and \(|\Phi(z)|\asymp|z|\).  For \(\Im z>0\),
every summand has nonpositive imaginary part.  If
\(\pi/4\le\arg z\le\pi/2\), the terms with \(2^n\asymp|z|\)
have negative imaginary part of magnitude at least a fixed multiple of
\(a_n\).  Thus the change in argument pushes the outer part of the
half-plane inside the prescribed sector.  Reflection proves the assertion
below the real axis; on the central part it follows from \(A=o(1)\).
Translating the half-plane sufficiently far to the right puts its whole
image in the domain of \(f\).

The bounded pullback \(F=f\circ\Phi\) is superexponentially small on a
fixed central sector, for example \(|\arg z|\le\pi/8\).  In a
half-disc of radius \(r\), the harmonic measure at a fixed positive
interior point of the central part of its semicircular boundary is at
least \(c/r\).  One can see this directly by scaling to the unit
half-disc: the Poisson kernel on a compact subarc away from its endpoints
is comparable to the distance of the rescaled axial point from the
straight boundary, namely a constant times \(1/r\).  Write
\(M\) for a bound for \(|F|\).  On that arc
\(\log|F|\le-Bc_1r+C_B\), and on the rest of the boundary
\(\log|F|\le\log M\).  The subharmonic maximum principle gives
\[
 \log|F(z_*)|\le\log M-c_2B+O_B(r^{-1}).
\]
Zeros cause no difficulty: use the subharmonic function
\(\max\{\log|F|,-N\}\) and then let \(N\to\infty\).
First let \(r\to\infty\), then \(B\to\infty\).  Hence \(F\) vanishes
at every fixed axial point.  The Identity Theorem gives \(F=0\), and
then \(f=0\), since \(\Phi\) is nonconstant and its image is open.
\end{proof}

We now describe the one-variable fields into which the selected paths
will be placed.  Begin with a positive variable \(x_1\to\infty\) and
the field \(B_0\) of convergent meromorphic Puiseux germs in \(1/x_1\).
At stage \(j\), the clock \(x_j\) belongs to \(B_{j-1}\), and for
\(j>1\) it satisfies \(x_{j-1}=o(x_j)\).  The first clock is \(x_1\),
or a positive constant multiple of it.  A germ in \(B_j\) has an
expansion
\begin{equation}\label{path:eq-clock-series}
 f\sim\sum_{\beta\in E}e^{\beta x_j}b_\beta,
 \qquad b_\beta\in B_{j-1},
\end{equation}
where the real set \(E\) is bounded above and finite above each cutoff.
The germ and the continued coefficients are holomorphic on admissible
sectors in the variable \(x_j\).  After a finite cutoff the remainder
has a strictly better exponential rate in \(\Re x_j\), on a possibly
smaller admissible sector.  This includes sector estimates for an empty
expansion.  The same definition is used for holomorphic families on
fixed ordinary-parameter polydiscs, with estimates on smaller polydiscs.
At the Puiseux level a parameter family has a common finite
ramification denominator and pole order, holomorphic coefficient
functions on the ordinary polydisc, and a convergent series locally
uniform there; bounds are taken on smaller polydiscs.

\begin{lemma}[Clock calculus and change of clock]\label{path:clock}
The preceding construction has the following properties, established
simultaneously by induction.
\begin{enumerate}[label=\textup{(\alph*)}]
\item Each \(B_j\) is a field, closed under differentiation in \(x_1\)
and conjugation.  A nonzero scalar has a leading equivalent
\begin{equation}\label{path:eq-leading-monomial}
 Cx_1^p\exp\!\left(\sum_{k=1}^j\beta_kx_k\right),
 \qquad C\ne0,
\end{equation}
uniformly on a suitable admissible sector.  Its real subfield is a Hardy
field: every nonzero real germ has a fixed eventual sign, and its
derivative either vanishes or has a fixed eventual sign.
\item If \(H\in B_j\) is positive and \(x_j=o(H)\), it is a valid
new sector variable.  Finitely many prescribed earlier-clock sectors
can be reached by taking a sufficiently narrow admissible \(H\)-sector.
If \(f\in B_j\) is positive, unbounded, and \(f=O(H)\), then on such
a sector its modulus is comparable to its value at the corresponding
positive real \(H=|H|\); it maps into any finitely prescribed
admissible \(f\)-sectors.  If \(f=o(H)\) on the real ray, then
\(\Re f=o(\Re H)\) throughout the narrowed sector.
\end{enumerate}
\end{lemma}
\begin{proof}
For \(B_0\), inversion and all assertions follow by convergent Puiseux
inversion.  A positive unbounded Puiseux germ is a positive constant
times a positive power of the real variable, with a relative error
tending to zero.  If \(f=O(H)\), the corresponding power is at most the
power for \(H\).  A strictly smaller power leaves a fixed angular
margin.  Equal powers require only an enlarged admissible loss.
Pullback of a loss under comparable positive power growth preserves
dyadic summability.

Assume both assertions at lower levels, and write \(x=x_j\).  Their
change-of-clock assertion permits evaluation of lower germs on
\(x\)-sectors and gives
\(\Re x_k=o(\Re x)\) for \(k<j\).  Consequently every fixed nonzero
lower scalar and its reciprocal have magnitude
\(\exp(o(\Re x))\).  The largest exponent in a nonempty expansion
therefore determines its leading term, separated from the rest by a
positive exponential gap.  Induction gives
Equation~\eqref{path:eq-leading-monomial}.  Sums and products are
calculated termwise: above any cutoff only finitely many pairs of
exponents occur.  To invert, factor the leading term and use a geometric
series in the exponentially small remainder.  A finite desired cutoff
uses only finitely many geometric terms.  A zero expansion is bounded
on the sector supplied by any one of its remainder estimates and
superexponentially small on the central sector.  Lemma~\ref{path:sector-uniqueness}
then makes the germ zero.  Thus the leading-term test is valid also
after cancellation.

Cauchy's formula on successively smaller sectors differentiates each
finite expansion and its remainder.  The radii of the Cauchy discs can
be chosen with logarithmic relative loss, after enlarging the angular
loss on neighboring shells.  Their reciprocal costs, and the
lower-field factor from differentiating the clock, can be absorbed in
an arbitrarily small loss of an exponential remainder rate.  The chain
rule and the induction hypothesis give closure under differentiation
in \(x_1\).  Reflection gives conjugation.  A real leading monomial
has a nonzero real leading constant, hence a fixed sign.  Applying this
to the derivative proves the Hardy-field assertion.  These steps use
only the lower-level change of clock, not the still-to-be-proved
change of clock in \(B_j\).

We prove that change next.  On every strip within a bounded distance
of the real \(x\)-ray, a nonzero lower scalar has ratio \(1+o(1)\) to
its value at a neighboring real point.  Indeed the earlier clocks have
size \(o(x)\), uniformly on fixed proportional discs by the induction
hypothesis, and Cauchy's estimate makes their derivatives \(o(1)\).
The logarithmic derivative of the factor \(x_1^p\) has the same
property.  The leading equivalent and a smaller sector give the
assertion for every nonzero lower scalar.  The assertion applies to a
nonzero lower-field derivative as well; division by \(dx/dx_1\)
expresses differentiation in \(x\) inside the lower field.

Write
\begin{equation}\label{path:eq-H-leading}
 H=e^{bx}h(1+E),\qquad h\in B_{j-1},\quad h>0,
 \quad E=O(e^{-\epsilon\Re x}).
\end{equation}
Since \(H\) is unbounded, \(b\ge0\).  First suppose \(b>0\).
On bounded strips,
\((\log H)'=b+o(1)\).  Analytic inversion along the arcs
\(H=\rho e^{i\theta}\), \(|\theta|\le\pi/2\), gives
\begin{equation}\label{path:eq-clock-inverse}
 x(\rho e^{i\theta})=x(\rho)+i\theta/b+o(1),\qquad
 \partial_\theta x=i/b+o(1),
\end{equation}
uniformly in \(\theta\).  For completeness, solve the inverse equation
first on a fixed small \(\theta\)-interval by the Analytic Inverse
Theorem, using the uniform lower bound \(|(\log H)'|\ge b/2\).
Continue over finitely many such intervals.  Integrating the inverse
derivative gives Equation~\eqref{path:eq-clock-inverse}, keeps the
solution in a bounded strip, and guarantees uniqueness on overlaps.
The resulting branches agree with the real inverse and hence agree
where their sector domains overlap.

If \(f=e^{ax}g(1+O(e^{-\epsilon'\Re x}))\), where \(g>0\) is lower
level, then unboundedness and \(f=O(H)\) give \(0\le a\le b\).
When \(a<b\), its argument tends uniformly to \(a\theta/b\), so
there is a fixed angular margin.  Its real size relative to \(H\)
decays exponentially in \(x\), which proves real-part separation even
after division by a logarithmic angular margin.  When \(a=b\), put
\(J=g/h\).  This is a positive bounded lower scalar, and
\[
 f/H=J(1+O(e^{-\epsilon_1\Re x})).
\]
The strip estimate implies, uniformly for \(|\theta|\le\pi/2\),
\[
 \frac{\Re f}{|H|}
   =O(J(x(\rho))\cos\theta)
     +O(|J'_x(x(\rho))|+e^{-\epsilon_2x(\rho)}).
\]
If \(J\to0\), its derivative is eventually negative unless it is
zero.  More precisely, put \(q=J'_x/J\), \(t=x(\rho)\), and
\(D=(\log H)'\).  Equation~\eqref{path:eq-clock-inverse} and the
strip estimate give
\[
 \frac{d}{d\theta}\arg J(x(\rho e^{i\theta}))
   =\Re\frac{q(x(\rho e^{i\theta}))}{D(x(\rho e^{i\theta}))}
   =\frac{q(t)}b(1+o(1)).
\]
Thus its phase changes toward the real axis for positive \(\theta\);
reflection handles negative \(\theta\).
If \(J\) has a positive limit, either direction of its vanishing
phase error is allowed.

The strip estimate also gives
\(J(x(\rho e^{i\theta}))=J(t)+O(|J'_x(t)|)\), and hence
\[
 \frac{|\Re f|}{\Re H}
 \le C J(t)+C\frac{|J'_x(t)|+e^{-\epsilon_2t}}{\omega_H(\rho)}.
\]
Choose the angular loss for \(H\) to dominate
\(|J'_x|+e^{-\epsilon_2x}\) by a factor tending to infinity.  This is
an admissible choice.  In fact \(J\) is eventually monotone with a
finite limit, so \(\int|J'_x|\,dx<\infty\).  Its suprema on bounded
\(x\)-intervals are comparable to its integrals there, by the strip
estimate for the derivative; a zero derivative needs no estimate.
Since \(\log H=bx+o(x)\) and its derivative tends to \(b\), these
intervals correspond to dyadic \(H\)-shells.  The preceding displayed
estimate now gives \(\Re f=o(\Re H)\) when \(J\to0\).  Prescribed
losses for the \(f\)-sector can also be pulled back: here
\(d\log f/d\log H\to1\), so a dyadic shell in either variable meets
only a bounded number of shells of comparable logarithmic width in
the other.  Their pulled-back suprema are summable.  Enlarging the
\(H\)-loss once more gives the sector mapping and modulus assertions.

It remains to consider \(b=0\).  Use \(h\) first as variable at the
previous level.  It is larger than \(x\), and \(\log h=o(x)\), because
the logarithm of a lower scalar is controlled by earlier clocks.
The lower change-of-clock assertion applies both to \(x\) and to the
leading lower factor of \(f\); \(f\) also has top exponent zero.
In the \(h\)-variable, the map from \(h\) to \(H\) in
Equation~\eqref{path:eq-H-leading} is a relative perturbation
\(O(e^{-\epsilon\Re x})\).  Take a wider and a narrower admissible
sector with margins at least a fixed multiple of
\((\log|h|)^{-2}\).  On the wider sector
\[
 \Re x\gtrsim |x|(\log|x|)^{-2},\qquad
 |x|\asymp x(|h|).
\]
Here are quantitative inverse estimates.  At a target point \(H_0\)
put \(r=|H_0|\), \(w=\log r\), \(X_0=x(r)\), and
\(\rho=cw^{-2}\).  Enlarge the losses on neighboring shells so that
\(D(H_0,4\rho r)\) lies in the wider sector.  The positive harmonic
function \(\Re x\) satisfies Harnack's inequality there, so on
\(D(H_0,2\rho r)\) the perturbation is bounded by
\[
 \eta=C\exp\{-c_1\epsilon X_0/(\log X_0)^2\}.
\]
Since \(w=o(X_0)\), one has \(\eta w^N\to0\) for each fixed
\(N\).  Rouch\'e's Theorem on a disc of radius \(\rho r\),
followed by Cauchy's estimate, gives a unique nearby inverse \(u\)
with
\[
 |u-H_0|=O(\eta r),\qquad
 \frac{dH}{dh}(u)=1+O(\eta/\rho).
\]
Uniqueness glues these inverse branches on overlaps.  For any positive
unbounded lower factor \(F=O(h)\), its sectorial logarithm is
\(O(w)\).  Cauchy's estimate therefore gives
\[
 |\log F(u)-\log F(H_0)|=O(\eta w^3)=o(w^{-2}).
\]
This proves modulus comparability and preserves every enlarged angular
margin.  If \(F=o(h)\), induction gives \(\Re F(u)=o(\Re u)\),
whereas \(\Re u/\Re H_0=1+O(\eta w^2)\).  The extra relative
exponential perturbation in \(f\) contributes at most \(O(\eta w^2)\)
after division by \(\Re H_0\), using a smaller common exponential
rate for the finitely many functions.  Thus real-part separation is
unchanged.  This completes the simultaneous induction.
\end{proof}

Only finitely many sector requirements are imposed at a time in
Lemma~\ref{path:clock}.  Different truncation orders may use different
losses and different initial radii.

\begin{lemma}[Sectorial limits of holomorphic families]
\label{path:family-limits}
Let \(F\) be a holomorphic family in a finite clock hierarchy,
on an ordinary polydisc \(D\).  Suppose that on the positive path
ray it converges to a finite holomorphic family \(F_\infty\),
locally uniformly on \(D\).  Then on every fixed smaller polydisc
\(D'\Subset D\), it converges uniformly to the same family on
some admissible sector of its last clock.  The conclusion remains
valid after passing to a faster admissible clock by
Lemma~\ref{path:clock}.  Finitely many such requirements can be
imposed on one sector.  In particular, a normalized limit which
is nonzero on a fixed compact contour stays uniformly nonzero
there, and finitely many composition arguments with limiting
images in prescribed open domains stay in those domains.
\end{lemma}
\begin{proof}
First note a uniform growth consequence of the parameter expansion
calculus.  A lower-level family, together with any fixed finite list
of its ordinary derivatives, is \(\exp(o(\Re x_j))\) on a suitable
sector of a faster clock \(x_j\), uniformly on a smaller ordinary
polydisc.  At the Puiseux level convergence gives a polynomial bound
in the first clock.  At a higher level take a finite cutoff below
the largest exponent.  There are only finitely many retained lower
coefficient families; their inductive bounds and the prescribed
remainder bound control the family by an exponential in its last
clock.  After passing to \(x_j\), Lemma~\ref{path:clock} makes the
real parts of all these earlier clocks \(o(\Re x_j)\); its modulus
comparability treats the polynomial factors.  Cauchy's formula on
one fixed smaller ordinary polydisc gives the same assertion for
any specified finite derivative list.  No nonzero normalization of
these families is required for this upper bound.

We now prove the finite-limit assertion by induction on the hierarchy.
At the Puiseux level, real-ray boundedness eliminates negative powers
of the small ramified variable.  The convergent series with holomorphic
coefficients then converges to its constant term also for complex
small argument, locally uniformly in the ordinary parameters.

At level \(j\), write the expansion as
\(\sum_\beta e^{\beta x_j}B_\beta\).  If a positive exponent has a
nonzero coefficient family, choose the largest such exponent and an
ordinary Taylor coefficient of \(B_\beta\) which is a nonzero scalar
germ.  Cauchy's formula makes the corresponding Taylor coefficient
of \(F\) bounded on the positive ray.  Its scalar expansion, however,
has a positive leading top exponent and a nonzero lower scalar
coefficient.  That coefficient and its reciprocal are subexponential
in \(x_j\) by Lemma~\ref{path:clock}, contradicting boundedness.
Zero coefficient families can be discarded; finiteness above each
cutoff ensures that the largest nonzero positive exponent exists if
there is one.  Thus every positive top coefficient vanishes.

Take the cutoff at zero and put \(B_0=0\) when that weight is absent.
The strict remainder gain gives
\[
 F=B_0+O(e^{-\epsilon\Re x_j})
\]
uniformly on a smaller ordinary polydisc and a suitable sector.  If
a truncation convention retains finitely many negative terms as well,
the preliminary uniform subexponential bound absorbs those terms
into this estimate with a smaller positive \(\epsilon\).
Consequently \(B_0\) has the same finite locally uniform limit on
the positive ray as \(F\).  By induction it has that limit on a
sector of the preceding clock.  Map the \(x_j\)-sector into this
sector by Lemma~\ref{path:clock}.  Modulus comparability ensures that
the preceding clock tends to infinity uniformly there.  This proves
convergence of \(B_0\), however slow its real-ray rate may be.
The displayed remainder tends uniformly to zero because
\(\Re x_j\gtrsim|x_j|/(\log|x_j|)^2\to\infty\).  The induction
is complete.

A further faster clock maps into the sector just obtained and sends
its tail to uniformly large moduli there, again by
Lemma~\ref{path:clock}.  Thus the limit persists.  A finite collection
of requirements uses a common admissible majorant of the finitely
many losses and a common sufficiently large radius.  Finally, the
modulus of a limiting unit has a positive minimum on a compact
contour.  Apply uniform convergence on a finite covering by smaller
ordinary polydiscs to make the error less than half this minimum.
For a composition argument choose a compact neighborhood of its
limiting image inside the prescribed domain and apply the same
uniform convergence.  These are finitely many requirements, as
claimed.
\end{proof}

\subsection{Regular systems and finite analytic operations}

We now collect the properties that must survive when another length is
introduced. Every nonzero analytic family must have a scalar
normalization; finite analytic operations must preserve controlled
domains; and the resulting expressions must still be realizable on
paths of the clock class. The definition below makes these requirements
an induction invariant. Preparation and root extraction will then be
finite operations within that invariant.

Fix a nonprincipal ultrafilter \(\mathcal U\) on the indices of the
given sequence.  Such a filter exists by the maximal principle applied
to the cofinite filter.  For every subset of indices, exactly one of
that subset and its complement lies in \(\mathcal U\).  Compactness
therefore gives limits in compact spaces, including the extended real
line.  All comparisons made at the distinguished sequence in the
construction below are in this sense.  Only a finite number of these
comparisons is used to construct any particular path.

\begin{definition}\label{path:def-regular}
A domain base is a directed family of connected real open domains in
auxiliary independent variables, each containing the distinguished
sequence for a set of indices in \(\mathcal U\).  A germ on the base is
an analytic function on one such domain, with its local complexification;
two functions are identified if they agree on a smaller domain.  Extra
ordinary variables range over complex polydiscs about specified standard
values.

A regular system \(\mathcal A\), with scalar field \(K\), has the
following properties.
\begin{enumerate}[label=\textup{(R\arabic*)}]
\item Every nonzero family \(F(q,\zeta)\) admits a nonzero scalar
\(k\in K\) such that \(F/k\) has a finite, nonzero holomorphic limit
as a germ in \(\zeta\), locally uniformly on a polydisc, along the
domain base.  The scalar reciprocal \(1/k\) is in the system.  For
real families \(k\) can be chosen real.  Every scalar comparison or
finite scalar limit made at the distinguished sequence holds along the
base after the appropriate restrictions.
\item For every positive scalar size \(a\) and integer \(n>0\), there
is a positive scalar \(\rho\) with \(\rho^n/a\) tending to a finite
positive constant.  Scalars and reciprocals have magnitude at most
\(\exp(C\max_i L_i^*)\), where \(L_i^*\) are the current centers.
For an empty tuple the bound is constant.
\item The system is closed under finite sums and products, reflection,
ordinary holomorphic composition at finite limiting arguments, inversion
after a nonvanishing normalized limit has been specified, and Cauchy
integration on fixed compact contours.  Integrals over fixed compact
segments are allowed by a finite cover of their holomorphic parameter
neighborhoods.  Domain restrictions for these operations are part of
the finite data of the expression.
\item Given a domain and finitely many normalization, comparison, and
limit requirements, there are families of real analytic paths eventually
in that domain preserving them.  The expressions used belong on each
path to a clock hierarchy of Lemma~\ref{path:clock}, with parameter
versions for any finite specified uses.  The path families have real
constant parameters in open sets and are analytic in those constants
where defined.  At every sufficiently large fixed path time, their
evaluation into the independent-variable domain is locally submersive.
The starting time may depend on the constants.
\end{enumerate}
Limits of families in \textup{(R1)} mean uniform limits on each fixed
smaller polydisc.  A nonzero holomorphic limit need not be nonzero at
its central point.  An inverse requires nonvanishing on the particular
polydisc or contour where that inverse is used.
\end{definition}

The expression ``finite analytic operation'' always includes its inputs
and the domain data needed for them.  This prevents an implicit request
to preserve infinitely many normalizing comparisons on a path.  Two
elementary consequences will be useful.

\begin{lemma}[Preparation on fixed contours]\label{path:prepared-contours}
Let \(F(q,\zeta,y)\) belong to a regular system.  Suppose that after
scalar normalization its limiting germ is regular of degree \(d\) in
the \(y\)-direction at \((\zeta,y)=(0,0)\).  Then its Weierstrass
polynomial, its unit, and division of any other family by that polynomial
are obtained by the operations in \textup{(R3)} on fixed contours and
fixed smaller parameter polydiscs.  Arbitrarily high Taylor coefficients
of these operations use the same denominator zero sets and the same
contours.
\end{lemma}
\begin{proof}
Choose a circle \(|y|=r\) on which the limiting germ is nonzero and
inside which it has precisely \(d\) zeros at \(\zeta=0\), counted
with multiplicity.  After shrinking the \(\zeta\)-polydisc and the
domain in the base, normalized \(F\) is uniformly nonzero on that
circle.  The power sums of its enclosed roots are
\[
 s_\ell(q,\zeta)=\frac1{2\pi i}\int_{|y|=r}
        y^\ell\frac{\partial_yF(q,\zeta,y)}{F(q,\zeta,y)}\,dy,
 \qquad 1\le\ell\le d.
\]
Newton's identities give the monic polynomial \(W\) from these power
sums, by finite polynomial operations.  The quotient \(F/W\) is a
holomorphic unit on a smaller disc by cancellation of the enclosed
zeros; its values there can also be expressed by a fixed-circle Cauchy
formula.  A useful explicit division formula for a holomorphic family
\(G\) is
\[
 G(y)=W(y)\frac1{2\pi i}\int_{|w|=r}
       \frac{G(w)}{W(w)(w-y)}\,dw
   +\frac1{2\pi i}\int_{|w|=r}
       \frac{G(w)}{W(w)}\frac{W(w)-W(y)}{w-y}\,dw.
\]
The second term is a polynomial in \(y\) of degree less than \(d\).
These formulas prove the asserted closure.  Differentiating or expanding
them introduces powers of the existing nonzero contour denominators,
not new denominators with new zero sets.  All Taylor orders are
holomorphic on the same smaller polydisc; their Cauchy constants may
depend on the order.
\end{proof}

\begin{lemma}[Roots in a regular scalar system]\label{path:roots}
Suppose \textup{(R1)--(R3)} hold, including divisibility of positive
sizes.  Then the complex scalar field is algebraically closed and its
real part is real closed.
\end{lemma}
\begin{proof}
We factor monic polynomials by induction on degree.  Translate the
variable to remove the coefficient of degree \(d-1\).  If all the
remaining coefficients vanish, the polynomial is a pure power and the
claim follows.  Otherwise write it as
\[
 y^d+\sum_{j=2}^d a_jy^{d-j}.
\]
Choose, from finitely many positive root-size representatives of the
nonzero \(a_j\), one of largest size, denoted \(\rho\).  Replacing
\(y\) by \(\rho v\) makes all coefficients bounded and makes at
least one non-leading coefficient have a nonzero limit.  The limiting
depressed polynomial has at least two distinct roots: a polynomial
with a single root is \((v-c)^d\), and its vanishing degree-\(d-1\)
coefficient forces \(c=0\).

Choose disjoint small fixed circles around the distinct limiting roots.
The contour power sums used in Lemma~\ref{path:prepared-contours}
give monic factors for the corresponding clusters.  Their coefficients
are in the system, and all their degrees are strictly smaller than
\(d\).  Induction factors each cluster completely.  Undo the scaling
and translation.  For a real polynomial, conjugation pairs the
nonreal factors and fixes every real-valued root.  The real scalar
field is ordered by eventual sign.  Every positive element has a real
square root, and every odd-degree real polynomial has a real root:
at the real points of a sufficiently small base domain these assertions
hold for the factorizations just constructed, and the finite sign and
conjugacy alternatives can be fixed on the domain base.  These two
properties characterize a real closed ordered field.
\end{proof}

\begin{lemma}[Detection of analytic identities]\label{path:detection}
Path families with the submersivity property in \textup{(R4)} detect
analytic germs.  Namely, if an analytic function vanishes on the tail
of every path in such a family, then it vanishes as a domain germ.
The same holds for a holomorphic family in extra ordinary variables.
\end{lemma}
\begin{proof}
Let \(V\) be an open parameter set for one family and use integer path
times \(N\).  At a fixed \(N\), evaluation is analytic on an open
subset \(V_N\subset V\).  On the subset where its differential is
surjective, the composition with a nonzero domain germ is a nonzero
real analytic function on every connected parameter neighborhood:
otherwise submersivity would give an open set of zeros in the domain.
The zero set of a nonzero real analytic function has Lebesgue measure
zero.  This last fact follows by induction on the number of variables:
in one variable zeros are discrete; in higher dimension separate the
points where a nonzero Taylor coefficient in the last variable is
nonzero and use the one-dimensional result and Fubini, with induction
on the exceptional coefficient-zero set.

Every parameter lies, for all sufficiently large integers, in a
submersive evaluation domain and in the zero set of the composition.
If the original germ were nonzero, countably many local parameter
charts for these integer evaluations would thus cover \(V\) by sets
of measure zero, a contradiction.  An open set of domain zeros makes
the germ zero by connectedness and the Analytic Identity Theorem.
For a family in extra variables, apply this argument to each of its
Taylor coefficients, or fix the extra variables in a countable dense
set and then use continuity.
\end{proof}

\subsection{Adjoining a reciprocal fastest length}

The next Lemma is the analytic specialization step.  It is stated with
its quantitative separation hypothesis, since that hypothesis rules out
uncontrolled cancellation against the lower field.

\begin{lemma}[One independent reciprocal]\label{path:adjoin-p}
Let \(\mathcal A\) be a regular system for a nonempty lower tuple
\(L'=(L'_1,\ldots,L'_r)\), with real scalar field \(K_{\R}\).
Assume that each \(L'_i\) is a positive unbounded element of
\(K_{\R}\), and put \(M'=1+\sum_{i=1}^rL'_i\).  Assume in addition
that, for every finite use of its path property, the lower paths can
be chosen with last clock \(x_{\rm prev}=O(M')\).  This is an
additional property of the length systems constructed below; it is
not part of the abstract conditions \textup{(R1)--(R4)}.
Let a new independent positive variable \(p\) have the distinguished limits
\[
 p\longrightarrow0,\qquad p\max L'\longrightarrow0,
 \qquad |p-c|\gtrsim p^2\quad(c\in K_{\R}).
\]
The constant in the last comparison may depend on \(c\).  Use domain
fibers in \(p\) which are open intervals cut out by finitely many
strict comparisons with lower scalars.  There is a regular system
\(\mathcal A[p]\) containing the lower system and \(p,p^{-1}\),
closed under the finite analytic operations, with the following size
bound for each nonzero scalar and its reciprocal:
\begin{equation}\label{path:eq-p-size}
 |a|+|a|^{-1}\le e^{C(1+\max L')}p^{-C}.
\end{equation}
Its finite path requirements can be realized by substituting for \(p\)
an expression in lower scalars, with an additional real constant varying
in an open interval, and preserving \(pM'\to0\).  On the resulting
lower paths one therefore has \(px_{\rm prev}\to0\).

For the empty lower tuple the same conclusion holds with ordinary
holomorphic germs and the choice \(p=\theta/x_1\), \(\theta>0\).
\end{lemma}
\begin{proof}
We give the chart construction, the normalization argument, and the
path argument separately.  Real closedness of the lower field is
available from Lemma~\ref{path:roots}.

\emph{Charts and their common refinements.}
For a real center \(c\in K_{\R}\), let \(s=|p-c|\); fix its sign on
the domain base.  If \(s\) is comparable to a positive lower scalar,
use a disc coordinate
\[
 u=(p-c)/R
\]
with finite limiting value, where \(R>0\) is in \(K_{\R}\).  If no
positive lower scalar has size comparable to \(s\), call the size
unmatched and use a ramified annular chart
\begin{equation}\label{path:eq-annulus}
 u=(s/R)^{1/n},\qquad v=(r/s)^{1/n},\qquad u,v\longrightarrow0,
\end{equation}
where \(R,r\) are lower positive scalars, \(r=0\) is allowed, and the
positive real branches are used.  In the latter chart permit scalar
prefactors \(ks^q\), with \(k\in K\), \(q\in\mathbb Q\).  Functions
are lower analytic families in these bounded coordinates, at their
specified standard values, times such prefactors.  Disc charts require
only lower scalar prefactors.

Here is the finite refinement rule.  Include zero and every center in
the finite data, and choose a center whose distance from \(p\) is
smallest in order.  If its distance is matched, a disc at that scale
expresses all other required coordinates by bounded analytic
substitutions, after normalizing their nonzero distances.  In the
unmatched case each nonzero difference of two centers is either
\(o(s)\) or much larger than \(s\): it cannot be comparable to \(s\),
since that would match \(s\) with a lower scalar.  Write
\[
 p-c'=(p-c)+(c-c')
\]
as the larger term times one plus the smaller ratio.  Choose annular
bounds \(r\ll s\ll R\) recording every one of these finitely many
ratios.  Rational powers are analytic in the factor tending to one,
and a common root denominator handles all ramification.  A required
bounded ratio of prefactors is treated in exactly this way: a ratio
\(ks^q\) has, for \(q\ne0\), a threshold at a lower root-size
representative of \(|k|^{-1/q}\).  Add these finitely many thresholds
to the bounds \(r,R\).  This proves the common-refinement rule,
including bounded analytic representation of every needed ratio.

\emph{Normalization in an annulus.}
On the relation \(uv=\lambda=(r/R)^{1/n}\), decompose a bounded
lower analytic family \(H(u,v,\zeta)\) as
\[
 H_+(u,\zeta)+H_-(v,\zeta),
\]
where the second series has only positive powers of \(v\).  This is
an actual convergent decomposition on a smaller bidisc.  For example,
substitute \((w,\lambda/w)\) on a fixed circle in the original
\(u\)-disc and take its nonnegative Laurent part by Cauchy projection;
the negative part is the positive series in \(v\).  Both belong to
the lower system, since \(\lambda\) is a bounded lower scalar and
the circle is fixed.  If \(r=0\), set \(v=0\) and use just the first
series.

Normalize a nonzero one of these families in the lower system, with
its series variable still ordinary.  Let \(N\) be the first degree
whose coefficient has nonzero standard holomorphic limit in
\(\zeta\).  The terms of degrees less than \(N\) have finitely many
lower scalar normalizers.  The tail beginning at degree \(N\),
factored by that power, is bounded with nonzero standard limit at the
origin.  Thus only finitely many candidate sizes occur.  Distinct
powers of \(s\) cannot have comparable sizes after lower scalar
multiplication: such a comparison, together with divisibility of lower
sizes, would match \(s\).  Choose the largest candidate size over
both Laurent parts.  It is unique up to terms with the identical
power, which are combined before this comparison.  All smaller
candidate ratios have zero limit, and the factored tails remain
bounded by the refinement rule.  Division by the largest size
therefore gives a bounded analytic representation with nonzero
holomorphic standard limit.  If both projected series vanish, the
original family vanishes.

\emph{Normalization in a disc.}
Translate the disc coordinate to have limit zero and normalize the
family \(H(u,\zeta)\) in the lower system before substituting \(p\).
If it is nonzero, choose a Taylor coefficient \(b(u)\) in \(\zeta\)
whose standard limit is a nonzero germ in \(u\).  Preparation and
division in \(u\), by Lemma~\ref{path:prepared-contours}, give
\begin{equation}\label{path:eq-disc-division}
 H(u,\zeta)=b(u)G(u,\zeta)+S(u,\zeta),
\end{equation}
where \(G\) is bounded and \(S\) is a polynomial in \(u\).  Also
\(b\) is a unit times a monic polynomial in \(u\).

The finite-dimensional \(K\)-span of the coefficient germs in
\(\zeta\) of \(S\) has a basis \(\psi_1,\ldots,\psi_h\) whose
standard limits are linearly independent over \(\C\), with each
\(\psi_j\) bounded.  To construct it, normalize a nonzero element,
choose one of its Taylor coefficients with nonzero limit, and normalize
that coefficient to be one.  Pass to the kernel of the corresponding
coefficient functional in the original span and repeat.  The dimension
drops by one at each step.  The resulting triangular Taylor-coefficient
matrix has nonzero limiting diagonal, proving the assertion.  Hence
\[
 S(u,\zeta)=\sum_{j=1}^h A_j(u)\psi_j(\zeta),
\]
where the \(A_j\) are polynomials over \(K\); its size after
specialization is the largest of their nonzero sizes.

Factor these finitely many polynomials, together with the polynomial
for \(b\), over \(K\).  Include the real parts of their roots, expressed
in the \(p\)-coordinate, among the centers of the refinement.  The
imaginary parts are lower scalars; their sizes are additional comparison
thresholds.  Every factor then has a bounded normalized analytic
representation, using either a disc or an annulus.  The largest size
among \(b,A_1,\ldots,A_h\) normalizes \(H\).  If an \(A_j\)-size
strictly dominates that of \(b\), independence of the standard
\(\psi_j\) prevents cancellation.  Otherwise \(b\) itself is a
Taylor coefficient of \(H\), so a normalized limit of \(H\) cannot
vanish identically.  The boundedness of the representation and Cauchy's
formula make this coefficient argument valid for holomorphic limits.

\emph{Bounds and analytic closure.}
All resulting sizes are products of lower scalar sizes and rational
powers of distances and scales.  Every nonzero distance from \(p\)
to a real lower center is at least a fixed multiple of \(p^2\).
Upper bounds for centers and scales, and bounds for their reciprocals,
are lower exponential bounds.  The preceding finite factorizations
therefore give Equation~\eqref{path:eq-p-size}.  Positive root-size
representatives follow from ramification and the lower representatives.
For a bounded scalar the chart formula gives its finite standard
limit; prefactor ratios are controlled by the recorded comparisons.
Nonzero real scalars have a fixed sign.  A complex normalizer for a real
family can be replaced by a real or imaginary part of comparable size,
and then by a positive representative if desired.

After common refinement, sums, products, ordinary composition, and
normalized unit inversion are the corresponding lower analytic
operations on bounded families.  Contour integrals are treated by
finitely many chart neighborhoods covering the contour.  The same
argument works for a compact segment.  This proves
\textup{(R1)--(R3)} for the extension.  Comparisons such as
\(|p-c|\gtrsim p^2\) can be expressed by interval endpoints in the
real lower field: factor the relevant real polynomial inequalities
using Lemma~\ref{path:roots}, order their roots, and choose the interval
containing the distinguished points.

\emph{Finite specialization and path families.}
Refine using all chart centers, all endpoints of the required interval,
and zero.  In a matched chart, with positive scale \(d\) and limiting
coordinate \(a\), choose
\begin{equation}\label{path:eq-matched-choice}
 p=c+ad+\theta d\eta,
\end{equation}
where \(\eta>0\) is a lower infinitesimal and \(\theta\) varies in a
small open interval of constants.  Its sign can be fixed to meet the
fiber constraints.  Every included endpoint has nonzero normalized
distance at its own recorded scale.  Choose \(\eta\) smaller than
the finitely many required relative margins.  Then all comparisons
persist and every required coordinate has its exact limiting standard
part.  If an endpoint were nearer than the chosen center-scale data,
it would already have been selected by the closest-center refinement;
thus this choice omits no cancellation relevant to the finite data.

In an unmatched chart choose a lower scale strictly between all
required lower and upper bounds, for instance their geometric mean
after taking the largest lower and smallest upper bound.  If the lower
bound is zero, multiply the upper bound by a positive lower
infinitesimal.  Set \(s\) equal to this scale times \(\theta\) in a
bounded positive open interval, and set \(p=c\pm s\) with its recorded
sign.  Thresholds required by bounded prefactors were already included,
so these substitutions preserve all finite normalization requirements.
They use only lower scalar and ordinary analytic operations.

Include among the finite requirements \(p\to0\) and \(pM'\to0\).
The latter uses only the already defined positive lower scalar
\(M'\): it is the requirement \(p/(1/M')\to0\), handled by the
same disc or annular refinement.  The lower system contains positive
infinitesimals, for example \(1/M'\).  Apply the lower path property,
with its assumed clock bound, to these and all other expressions and
comparisons used in the chosen substitutions.  Then
\(px_{\rm prev}=O(pM')\to0\); no comparison involving a clock
chosen only later is imposed on the domain base.
Formula~\eqref{path:eq-matched-choice}, or
its unmatched counterpart, has nonzero derivative in the new constant
\(\theta\).  Together with submersivity of the lower evaluation,
the triangular differential gives submersivity into the extended
domain.  Interval fibers over connected lower domains are connected
after ordering their endpoints; they remain the domain base under
these restrictions.  This proves \textup{(R4)}.  In the empty case,
  \(p=\theta/x_1\) gives the stated meromorphic Puiseux families directly.
\end{proof}

\subsection{Construction for several lengths}

We now construct the regular system of
Definition~\ref{path:def-regular} for the original tuple of length
variables. The induction separates lengths of comparable size from
strictly slower differences, applies the reciprocal-length construction
when necessary, and preserves the common coefficient domains. Its
output will leave only the ordinary-variable sign conditions to solve
in the final subsection.

For example, the admissible equation
\[
 e^{-L_1}-L_2e^{-L_2}=0
\]
has the path \(L_2=e^u\), \(L_1=e^u-u\).
The lengths have ratio tending to one, but their difference is the
unbounded slower quantity \(-u=-\log L_2\). Thus grouping lengths by
their limiting ratios must retain the magnitudes of unbounded slower
differences as new centers.

\begin{proposition}[Centers and regular systems]\label{path:centers}
For a distinguished sequence of \(m\) lengths tending to infinity,
there are centers \(L_1^*,\ldots,L_m^*\), differing from the sequence
values by bounded additive amounts, and a regular system on a domain
base with the following additional properties.
\begin{enumerate}[label=\textup{(\alph*)}]
\item The system contains the centers, \(e^{bL_i^*}\) for every real
\(b\), and every admissible family evaluated at these centers.
\item Every finite expression built with the allowed analytic
operations has coefficient formulas on a common connected domain of
the base and a common ordinary-parameter neighborhood.  The same
formulas remain valid to all orders on paths preserving its finite
preparation data.  Remainder sectors and thresholds may depend on
order.
\item On the path families, the grouped centers have the form
\(c_iX+D_i^*\), \(c_i>0\), where the positive magnitudes of the
unbounded \(D_i^*\) are slower centers.  All required quadratic
holomorphy domains are reached after narrowing the clock sectors.
\item For a nonempty tuple, its path families can be chosen so that
the last clock is the selected fastest center \(X\).  Every center
is \(O(X)\), and one of the centers equals \(X\).  In particular
the last clock is comparable to \(\max_iL_i^*\), independently
of which finite expression requirements are imposed.
\end{enumerate}
\end{proposition}
\begin{proof}
We induct on \(m\), establishing all regularity and path assertions
together, including the last-clock invariant in \textup{(d)}.
The empty tuple has ordinary holomorphic germs, constant
scalars, and a point as its independent-variable domain.  Its
normalization and contour assertions are immediate; the path assertion
has zero-dimensional target and is vacuous.

\emph{Splitting the fastest block.}
Choose a fastest length \(X\), so that every ratio \(L_i/X\) has
a finite ultrafilter limit.  This is possible by taking a maximal
one among the finitely many lengths.  For the indices with positive
limit write
\[
 L_i=c_iX+D_i,\qquad c_i>0,\quad D_i=o(X).
\]
Ignore bounded \(D_i\) for now, replacing them by zero in the centers.
For each unbounded difference take its absolute value and record its
eventual sign.  These magnitudes and the lengths with ratio zero to
\(X\) form a lower tuple of at most \(m-1\) entries.  Indeed a
fastest block of size \(k\) contributes at most \(k-1\) differences,
and the other block has \(m-k\) lengths.  Apply the induction
hypothesis to that tuple and denote the resulting positive centers by
\(L'\).  Reinsert the recorded signs in the differences \(D_i^*\).

If some real lower scalar approximates the original \(X\) to bounded
absolute error, replace \(X\) by it.  Its positivity, its divergence,
and \(\max L'=o(X)\) are finite lower-system requirements.  Otherwise
adjoin \(p=1/X\) as an independent variable.  For every real lower
scalar \(c\),
\begin{equation}\label{path:eq-separation}
 |p-c|\gtrsim p^2.
\end{equation}
In fact failure, or even a bounded ultrafilter value of
\(|p-c|/p^2\), would give \(c/p\to1\) and
\[
 X-1/c=\frac{c-p}{pc}=O(1),
\]
contrary to the case distinction.  Thus Lemma~\ref{path:adjoin-p}
applies.  In both cases call the resulting coefficient system
\(\mathcal A'\).  Every nonzero scalar of \(\mathcal A'\), and its
reciprocal, is \(e^{o(X)}\), by
Equation~\eqref{path:eq-p-size}, \(\max L'=o(X)\), and
\(\log X=o(X)\).  The final centers in the fastest block are
\(c_iX+D_i^*\), including \(X\) itself.  All errors relative to the
distinguished original tuple are bounded.

\emph{Initial grouped expansions and their domains.}
Adjoin \(e^{\alpha X}\), \(\alpha\in\R\), and the initial
admissible functions evaluated at the centers.  Their fastest-block
expansions have the form
\begin{equation}\label{path:eq-grouped}
 \sum_d e^{-dX}B_d,\qquad B_d\in\mathcal A',
\end{equation}
with locally finite weights and strict exponential remainder gains.
To obtain the coefficient, expand a term with fastest-block indices
\(d_i\) as
\[
 \prod_i e^{-d_i(c_iX+D_i^*)}(c_iX+D_i^*)^{k_i}.
\]
Its top exponential is \(e^{-(\sum_i c_id_i)X}\).  Its remaining
factors are polynomials in \(X,D_i^*\), exponentials of the signed
lower centers, and recursively extracted initial coefficients in the
slower lengths.  All are in \(\mathcal A'\).  Definition~\ref{path:def-admissible}
supplies both the grouped remainder and the common analytic domains
of every coefficient formula.

The inductive lower construction supplies
\(x_{\rm prev}=O(1+\sum_iL'_i)\), and contains the positive lower
centers as real scalars.  Thus it meets the additional hypotheses of
Lemma~\ref{path:adjoin-p}.  In the independent case that specialization
preserves \(p(1+\sum_iL'_i)\to0\); hence
\(x_{\rm prev}=o(X)\), where \(X=1/p\).  In the dependent case
\(X\) is a lower scalar, and the already required comparison
\((1+\sum_iL'_i)/X\to0\) can be imposed using the lower path
property; it again gives \(x_{\rm prev}=o(X)\).  Thus in either
case \(X\) belongs to the lower path field and is an admissible
next clock by Lemma~\ref{path:clock}.  Declare it to be the last
clock of the new hierarchy.  All fastest-block centers are
\(c_iX+o(X)\), the other centers are \(o(X)\), and one center
is \(X\), proving \textup{(d)}.  When the lower tuple is empty,
the independent choice \(p=\theta/x_1\) gives the first clock
\(X=x_1/\theta\).  Subsequent finite analytic operations and domain
requirements introduce no further clocks.
The slower centers and the positive magnitudes of the differences
satisfy
\[
 \Re L'=o(\Re X),\qquad |L'|=o(|X|)
\]
on a narrowed sector.  Consequently
\(\Re(c_iX+D_i^*)\asymp\Re X\) and
\(|c_iX+D_i^*|=O(|X|)\).  The admissible loss gives
\(\Re X\gtrsim |X|/(\log|X|)^2\), so for every fixed \(c>0\)
and sufficiently large \(|X|\),
\[
 |\Im(c_iX+D_i^*)|
 \le C|X|<c\bigl(\Re(c_iX+D_i^*)\bigr)^2.
\]
Bounded ordinary shifts satisfy the same estimate.  This verifies
evaluation in every required quadratic region.

This initial substitution assertion is unconditioned: it holds for
every recursively extracted initial coefficient on every path just
described, without prescribing its leading nonzero size.  At a slower
grouping level one repeats the same expansion, using its positive
ratios and separated smaller real parts.  Only initial coefficient
expansions, clock powers, and clock exponentials occur.  This observation
is essential because infinitely many initial coefficients can occur
in the all-order formulas for one finite expression.

\emph{Closure for prepared finite expressions.}
Consider an expression tree formed by the allowed operations.  At
each inverse node record the input's first nonzero coefficient, a lower
scalar normalizing it, and a fixed disc or contour on which its
normalized limit is nonzero.  At each composition node record the
finite limiting arguments and a pair of nested polydiscs inside the
analyticity domain of the outer function.  At each contour or segment
integral record a finite cover of the fixed integration set.  At a
preparation or division node use the fixed-circle formulas of
Lemma~\ref{path:prepared-contours}.  Call these finitely many records
the preparation data of the expression.  They are constructed from
the bottom of the expression tree upward; no inverse is taken before
its input has been classified.

Sums and products preserve Equation~\eqref{path:eq-grouped}.  At an
inverse node, factor the first nonzero exponential term and invert
its coefficient on the recorded unit domain.  A geometric expansion
then gives all coefficients and strict remainder gains.  At a
composition node a bounded family has no nonzero growing exponential
term, because at an ordinary argument where the first normalized
coefficient is nonzero such a term would be unbounded.  Its full
weight-zero coefficient is bounded too, by the same lower normalization
argument.  Expand the outer holomorphic function at that full
weight-zero coefficient.  The remaining argument is exponentially
small; a finite Taylor polynomial suffices at each cutoff, and the
Cauchy remainder on the fixed nested polydiscs has a strictly better
exponential order.  If there is no weight-zero term, use zero as
the center.  Contour and segment integration preserve the estimates
on the finite covering patches.

These rules describe every coefficient by a finite formula.  At
arbitrary coefficient order, that formula can involve derivatives of
arbitrarily high order, but only powers and derivatives of the
\emph{same} prepared denominators, and evaluations on the \emph{same}
prepared domains.  A derivative of a function holomorphic on a fixed
polydisc is holomorphic on that polydisc; all its Cauchy bounds are
available on one fixed smaller polydisc, with order-dependent constants.
There are only finitely many nodes in the expression tree.  Choosing
nested neighborhoods for these nodes once therefore gives a common
ordinary-parameter domain for all coefficients.  The common initial
coefficient domains and the finitely many recorded restrictions give
a common connected domain in the base as well.  Increasing an
asymptotic threshold does not change these analytic formulas.

\emph{The path assertion for those same formulas.}
Put the finitely many lower families in the preparation data into
the charts of Lemma~\ref{path:adjoin-p}, where needed.  Specialize
\(p\) there and impose all their finite domain and normalization
requirements on lower paths.  Such chart substitutions are regular
lower operations.  In the empty lower case they are meromorphic
ramified functions of \(1/x_1\).
Lemma~\ref{path:family-limits} gives sectorial convergence of each
of the finitely many bounded normalized lower families in the
preparation data, locally uniformly on the chosen ordinary
polydiscs.  The compact-contour and composition conclusions of
that Lemma keep the recorded units nonzero and the recorded
arguments inside their domains on one common sector.  Thus the
contour estimates are valid there.

Every unconditioned initial coefficient, including ones not listed
in the preparation data, has a clock expansion by the earlier initial
substitution argument.  Every further coefficient is made from these
by the finite-expression rules just proved: sums, products, ordinary
derivatives, evaluations and unit inversions on the prepared domains,
and fixed integrals.  The lower path calculus therefore gives those
same coefficient formulas and remainder estimates at all orders.
Coefficients may unexpectedly vanish on a particular path; their
formulas and the remainder estimates still hold.  This does not
require the leading sizes of infinitely many coefficients to be
preserved.  Sectors and thresholds can be narrowed separately at
each order.

\emph{Normalization and the zero-series alternative.}
If an expression has a first nonzero coefficient in
Equation~\eqref{path:eq-grouped}, normalize that coefficient in
\(\mathcal A'\) and multiply its scalar normalizer by the corresponding
exponential.  Every lower nonzero scalar and its reciprocal is
subexponential in \(X\), so the exponential gap makes the remainder
negligible after this normalization, locally uniformly on the ordinary
polydisc.  This supplies \textup{(R1)}.  The size is an exponential
times a lower scalar size; it is divisible up to a positive constant,
and it and its reciprocal obey the required exponential growth bound.
Scalar signs and bounded scalar limits follow from the same leading-term
calculation.  Real normalizers follow by reflection and choosing a
real component of nonzero leading size.

If every coefficient is zero as a lower germ, each coefficient
vanishes on its entire common connected domain by the Identity Theorem.
This uses common analyticity, not a common remainder threshold.  On
every path preserving the finite preparation data, the expression
therefore has zero clock expansion.  One remainder estimate gives
boundedness on a sector, while estimates of arbitrary order give
superexponential decay on the central subsector.  Lemma~\ref{path:sector-uniqueness}
makes the expression zero on every such path.  The path families
are submersive by Lemma~\ref{path:adjoin-p} and the induction
hypothesis; Lemma~\ref{path:detection} consequently makes the expression
zero as a domain germ.  This proves the alternative used to classify
inputs before further inversions.

For finitely many expressions, preserve the finite coefficient
normalizations and comparisons just identified.  The preceding path
argument proves \textup{(R4)} and the claimed clock membership, while
the closure rules prove \textup{(R3)}.  Lemma~\ref{path:roots} then
supplies scalar roots at the new level for the next induction step.
The exponential growth bounds have maximum center comparable to \(X\),
so \textup{(R2)} also holds.  The induction is complete.
\end{proof}

\subsection{Eliminating ordinary variables and preserving their limits}

We record the algebraic sign test used to eliminate one ordinary
variable. Its mechanism is the classical Hermite trace form and sign
determination from Tarski queries; see
\cite[arXiv version, Theorem~7 and Proposition~11]{PerrucciRoy2017}.
We include the finite calculation and check its use over our regular
scalar field. The multivariable path construction remains the one
proved above.

\begin{lemma}[Univariate sign elimination]\label{path:sign-elimination}
For a finite list of real polynomials \(P_i(y)\) of bounded degrees,
existence of a real \(y\) realizing any specified Boolean combination
of their signs is determined by finitely many polynomial sign tests
in their coefficients.  The same tests are valid for polynomial
coefficients that are real scalar germs in a regular system.
\end{lemma}
\begin{proof}
First separate the finitely many alternatives for degrees and
identically zero polynomials.  Constant signs are already coefficient
signs.  Let \(D\) be the product of the nonconstant polynomials.
On the real line, every sign pattern is attained at a root of \(D\),
in a bounded complementary interval, or in one of the two unbounded
intervals.  Rolle's Theorem gives a root of \(D'\) in every bounded
interval between consecutive distinct roots of \(D\); none of the
\(P_i\) vanishes inside such an interval, so all its signs are constant.
The two unbounded intervals are tested by degree and leading-coefficient
signs.  It therefore suffices to count sign types at the distinct real
roots of \(Q=D\) and, when nonconstant, \(Q=D'\).

For \(h_i\in\{0,1,2\}\), put \(P=\prod_iP_i^{h_i}\), with the
zeroth power interpreted as one, and consider
\begin{equation}\label{path:eq-trace-form}
 \mathcal H_{Q,P}(v,w)
 =\operatorname{Tr}_{\R[y]/(Q)}(Pvw).
\end{equation}
The right-hand side is the trace of multiplication by \(Pvw\).
Its signature is
\[
 T_h=\sum_{a:\,Q(a)=0,\ a\in\R}^{\text{distinct}}
             \prod_i\operatorname{sign}(P_i(a))^{h_i}.
\]
To verify this even for repeated roots, decompose the algebra into its
local factors.  At a real root of multiplicity \(m\), trace is
\(m\) times evaluation; nilpotent directions are in the radical,
and the remaining one-dimensional form has sign
\(\operatorname{sign}P(a)\).  At a nonreal conjugate pair, the value
coordinates form one complex line viewed as a real plane.  If
\(P(a)\ne0\), the form is a nonzero multiple of
\(\Re(P(a)vw)\) and has one positive and one negative direction;
if \(P(a)=0\), it is zero.  Its signature is therefore zero.

For a sign \(s\in\{-1,0,1\}\), its three indicator polynomials are
\[
 \mathbf 1_{s=0}=1-s^2,\qquad
 \mathbf 1_{s=1}=\tfrac12(s^2+s),\qquad
 \mathbf 1_{s=-1}=\tfrac12(s^2-s).
\]
Multiplying these identities shows that the number of roots of any
specified sign type is a rational linear combination of the finitely
many \(T_h\).  Its positivity tests existence of that type.

In a monomial basis the matrix of Equation~\eqref{path:eq-trace-form}
is obtained by rational operations on polynomial coefficients, after
the leading coefficient of \(Q\) is known to be nonzero.  Its signature
is determined by finitely many sign tests: pivot on a nonzero diagonal
entry and pass to its Schur complement; if the diagonal is zero but
an off-diagonal entry is nonzero, replace a basis vector by the sum
of the corresponding two vectors to obtain a nonzero diagonal pivot.
A zero matrix ends the procedure.  This finite algorithm, with its
finitely many pivot alternatives, uses only rational functions of
the coefficients.  Clearing denominators with their recorded signs
gives polynomial sign tests.

For scalar germs, factor the polynomials using Lemma~\ref{path:roots},
fix the order and real or conjugate status of the finitely many roots,
and restrict the domain base accordingly.  At every real point of
that domain the preceding real-polynomial argument applies, with the
same sign alternatives.  Thus the same finite tests hold as germ
statements; no transfer principle or additional quantifier-elimination
theorem is needed.
\end{proof}

\begin{proof}[Proof of Theorem~\ref{thm:path-selection}]
Apply Proposition~\ref{path:centers} to the given lengths.  Add the
bounded differences between the original lengths and their centers
to the ordinary variables.  Their ultrafilter limits are finite.
For an initial admissible family, these bounded shifts preserve
admissibility: a bounded complex shift maps a smaller quadratic
region into the original one; exponential-polynomial terms transform
by ordinary analytic operations; and a smaller common fixed-width
band remains inside the original common bands.  Thus its shifted
values belong to the regular system of the Proposition.

There are only finitely many complete sign patterns for the functions
in a Boolean formula.  Choose one realized on a set in \(\mathcal U\).
Translate all ordinary limiting values to zero.  Normalize every
nonzero family by a real scalar, recording its fixed sign.  A family
with nonzero limiting value at zero is a unit of fixed sign on a
smaller domain; an identically zero family has known sign.  For the
remaining finite list choose one real direction in which every
nonzero limiting germ is regular.  Such a direction exists: the
lowest nonzero homogeneous Taylor polynomial of each germ excludes
a proper real algebraic subset, and a finite union of those subsets
does not exhaust the space of directions.  Complete this direction
to a fixed real linear coordinate change.

Use the last coordinate \(y\) in that direction.  By
Lemma~\ref{path:prepared-contours}, each remaining condition is the
sign condition of a monic polynomial in \(y\), multiplied by a unit
of fixed sign.  At the limit of the remaining ordinary variables its
polynomial tends to a pure power \(y^d\).  All its roots therefore
tend to zero: if its non-leading coefficients tend to zero, the
elementary root bound applied outside any fixed disc excludes roots
there for sufficiently small coefficients.  Lemma~\ref{path:sign-elimination}
replaces existence of the designated signs by finitely many coefficient
sign tests in one fewer ordinary variable.  Choose the realizable
finite alternative at the distinguished sequence and repeat.  Each
step uses normalization and a fixed generic direction for only its
finite list; it is an operation of the regular system with finite
preparation data.

Once no ordinary variables remain, all conditions are scalar
comparisons on the domain base.  Lift the choices in the reverse
order.  The relevant polynomial roots belong to the scalar field by
Lemma~\ref{path:roots}.  A realizable sign pattern can be attained at
a real root, at the midpoint of two consecutive real roots, or just
beyond an extreme root.  For the last choice add or subtract any
sufficiently small positive scalar infinitesimal.  If there are no
real roots, take zero; if necessary a positive or negative infinitesimal
also gives the required constant interval pattern.  All roots tend
to zero, so every such choice tends to zero.  Finitely many required
neighborhood bounds on units can be imposed when choosing the
infinitesimal.  The previously chosen ordinary variables already have
their prescribed limits.  Repeating the lifting realizes every
ordinary variable as a scalar-system element, with all signs and
limits valid on a domain of the base.

Now choose a path using the finite uses of these scalar operations,
the finite normalization and unit requirements, and the additional
finite admissible list in the statement.  This is precisely the
finite-data path assertion of Proposition~\ref{path:centers}.
Undo the fixed linear coordinate changes and restore the bounded
shifts.  The \emph{actual} coordinates are sums and linear combinations
of scalar-system elements on the path, so they, and each specified
additional function, lie in a real Hardy field by
Lemma~\ref{path:clock}.  Their derivatives either vanish or have
fixed eventual sign.  Hence these actual coordinates and functions
are eventually monotone or constant.  All their required limits were
preserved in the construction.  This proves the Theorem.
\end{proof}
